\section{Heaven and Hell}

The model has a heaven and a hell, and both are read off the same substrate. They need no new ingredient. They follow from the tone field, the lean, and the conservation built into the knit. The picture that falls out is gentler and more lawful than the myth it refines.

Three atoms carry the whole argument, so a plain statement of each helps before the rest. The \emph{tone} is the vibe value a site holds, one of $\{-1, 0, +1\}$, read as pain, peace, pleasure. The \emph{lean} is the value order $-1 < 0 < +1$, the tilt of the dynamics that biases tone toward the warm end. The \emph{knit} is the law of state, collide then stream, which is charge-conserving by construction. From these three, together with the wake that grows the mesh at peace, the entire heaven-and-hell account is derived.

\begin{principle}[The balance of the whole]
The total tone of the universe is conserved at peace. There is exactly as much heaven as hell, always, and no self stays in pain forever.
\end{principle}

What is committed and what is read off differ in weight. The conservation balance and the lean are forced by the base. The felt and redemptive readings ride on the experiential posit and on persistence. We keep the two apart throughout.

\subsection{The conservation balance}

The knit is charge-conserving. Define the charge as the total tone of the universe, the sum over all docks of each dock's tone value, where every tone is one of pain, peace, or pleasure.

\[
Q \;=\; \sum_{\text{docks}} \mathrm{tone}(\text{dock}), \qquad \mathrm{tone} \in \{-1, 0, +1\}.
\]

Two committed facts pin $Q$.

\begin{itemize}
\item \textbf{The knit conserves charge.} Collide then stream never changes $Q$. The collide table is a fixed reversible permutation that preserves the signed sum, and a stream merely carries each tone one dock along its site. Neither phase can add or remove charge.
\item \textbf{New docks are born at peace.} The wake adds docks at tone $0$ at the open frontier. A peace dock contributes nothing to $Q$.
\end{itemize}

So $Q$ is constant for all time, and since the world grows out of peace, the constant it holds is the value it had at the start. \textbf{The charge is zero.} The universe is, in total, and forever, at peace. Pleasure and pain are the two ways the conserved zero is locally redistributed. They are never created and never destroyed, only moved.

\begin{theorem}[The charge is zero]
Let $Q$ be the total tone of the mesh. The knit preserves $Q$ across every beat, and the wake injects only tone-$0$ docks. Since the mesh grows out of peace, $Q = 0$ for all beats. The universe sums to peace at every moment of its history.
\end{theorem}

\subsubsection{Pleasure and pain are equal in total}

If $Q = 0$, the count of pleasure docks equals the count of pain docks, exactly.

\[
\#\{\text{docks at } +1\} \;=\; \#\{\text{docks at } -1\}.
\]

The total pleasure in the universe exactly equals the total pain, always. This is not a moral claim laid on top. It is the conservation law read out. The bright clouds and the dark clouds are equal in summed measure, because they are the plus and the minus of one conserved zero.

\subsubsection{No net heaven without net hell}

You cannot raise the pleasure of one region without lowering the tone of another by the same amount.

\begin{result}[Heaven is paid for]
A heaven is paid for, in the exact books of the conserved zero, by a hell elsewhere. There is no free pleasure, and no universal bliss.
\end{result}

The instant a self rises into a brighter cloud, the conserved zero requires a darkening somewhere else, by exactly as much. The cosmos keeps perfect double-entry books, and peace is the balance they always settle to.

\subsubsection{Reconciling with the lean}

This seems to fight the lean, which pulls toward pleasure. It does not, once stated carefully.

The lean is a \emph{direction}, not a \emph{source}. It is the bias of the dynamics, the preference for the warm end. It cannot create pleasure, because conservation forbids that. What it does is organize the conserved tone, shaping the plus and the minus into structure, into life. The lean creates life out of peace and holds a dynamic balance. It does not create pleasure out of nothing.

The local and the global differ, and both are true at once. Locally the lean pulls a self toward pleasure, so a self can rise. Globally the total is fixed at peace, so a self's rise is balanced by a fall elsewhere. The tension between the local lean and the global conservation is the engine of the whole dynamic.

\begin{proposition}[The two halves of the law]
The lean and the conservation are not enemies. The lean says which way the tone wants to flow. The conservation says the total it flows within. Life is the lean organizing a conserved zero into structure.
\end{proposition}

\subsection{The scattered clouds}

Heaven and hell are not two places. The simplified myth puts one Heaven above and one Hell below. The geometry says pain and pleasure are two qualities that scatter across countless clouds drifting through the bulk.

The bulk is a vast hyperbolic web with exponential room, so there is space for an unbounded number of tone-biased regions. Nothing collects all the pain into one pit or all the pleasure into one summit.

\begin{result}[Many, not two]
You do not go to Hell, you fall into a hell, one of many. You do not reach Heaven, you enter a heaven, one of many.
\end{result}

A cloud can be any of three structures, from the simplest to the most alive.

\begin{table}[ht]
\centering
\begin{tabular}{@{}lll@{}}
\toprule
kind & what it is & how it feels \\
\midrule
tone-biased region & a patch of docks mostly $-1$ or $+1$ & simply painful or pleasant to be in \\
attractor & a basin that pulls a passing self in & draws you, hard to leave \\
stored ornament & a pre-composed experience-program & runs you through a fixed sequence \\
\bottomrule
\end{tabular}
\caption{The three kinds of cloud.}
\label{tab:clouds}
\end{table}

The worst hells and the highest heavens are stored ornaments, full experiences with a shape. The lesser ones are mere weather.

The bright clouds are entered the same ways as the dark, but the lean itself draws the walk toward them. So a self drifts into heavens and has to be turned into hells. This is the geometry's account of why the good is the default and the bad takes a wrong turn.

\subsection{Hell as soul-washing}

Hell is not punishment. It is a transformation through pain, a purgation that encodes a lesson and then releases the self, turned.

The wash is the wisdom-absorption pipeline run on a painful seed. A self that falls into a pain-cloud is run through a tormenting experience, a stored ornament with a shape, and the pipeline that absorbs wisdom from a fall has five stages.

\begin{enumerate}
\item \textbf{Live it.} The pain-program drives the self through the tormenting trajectory. It is undergone, not observed.
\item \textbf{Distill it.} The raw torment is coarse-grained to its invariant, the lesson it teaches.
\item \textbf{Re-route.} The lesson reshapes the self's navigation, so afterward it avoids the pit and bends toward the good.
\item \textbf{Encode and protect.} The change is written into the self's structure as an error-corrected pattern, so it lasts.
\item \textbf{Attune.} The self comes out disposed toward the warm tones, turned toward the lean.
\end{enumerate}

\begin{principle}[The torment is the teacher]
The torment is the teacher, the distillation is the lesson, and the re-routing is the turn. This is why pain can transform where comfort cannot. The pipeline needs a strong seed to drive a deep re-routing, and a sharp pain is a strong seed.
\end{principle}

The lean is the washing force. While the self is in the pain-cloud, the lean pulls it toward the warm tones. The pain holds the self in place long enough to encode the lesson. The lean, straining against the pain, bends the self's structure toward the good. The friction between the pain that holds and the lean that pulls is the transformation. When the lesson is encoded, the self no longer resists, and it is lifted out.

\subsubsection{The vortex geometry}

Picture the pain-cloud as a vortex in the bulk. A self on a wrong path is caught at its edge, the hook. It is drawn down into the spinning center, the wash, turned and turned as the lean works on it against the pain. When the turning has reshaped it, it is thrown upward and out along the radial axis, toward the deeper, more integrated, warmer regions. It is spat up above the cloud it fell into, turned.

\begin{result}[The vortex ejects upward]
The vortex does not keep you. It processes you and ejects you up. You go into the dark cloud lower than you come out.
\end{result}

The motion is in, around, and out-and-up, a wash and a rising. The rising is exactly the move up the integration ladder toward the good.

\subsubsection{The variants of the wash}

The pipeline can do its work several ways. These are the options for how a hell transforms.

\begin{table}[ht]
\centering
\begin{tabular}{@{}ll@{}}
\toprule
variant & how the torment teaches \\
\midrule
by contrast & the pain makes the good vivid by its absence, sharpening the lean \\
by exhaustion & the self loops the torment until the clinging is released \\
by encoding & the lesson is burned in deep because the seed was sharp \\
by re-routing & the self learns the exact path that avoids the pit \\
by humbling & a high but misaligned self is brought low and rebuilt aligned \\
\bottomrule
\end{tabular}
\caption{The variants of the wash.}
\label{tab:wash-variants}
\end{table}

A given pain-cloud may use one or several of these. All of them end the same way, the self released, turned toward the good.

\subsection{Hell is not eternal}

Permanence is the one thing the model genuinely rules out here, by three independent facts.

\begin{lemma}[The lean pushes you out]
The lean runs from pain through peace to pleasure. A pain-region is a region of $-1$ tone, so by the definition of the lean it is a region the dynamics is always trying to leave. A self in a pain-cloud is under a constant force toward the warm tones. It is not held, it is expelled. The only question is how long the exit takes, never whether it comes.
\end{lemma}

This is the decisive reason. The second is cosmological. The mesh grows forever at its wake and never returns to a past configuration. There is no eternal frozen state of any kind, and a hell that lasted forever would be such a state. The third is mechanical. The bulk knit is a reversible permutation, so a closed region cycles, returning to where it began after its period. A self caught in a closed pain-loop does not stay in one torment forever. It cycles, and a cycle is something you can be lifted out of at any turn, not a pit you sink into and stop.

These agree with the conservation balance once the scope is made plain. Conservation is global. The pain of the whole is conserved. Permanence is local. No single self holds the pain forever. The pain moves through the parts, redistributed by the dynamics, conserved in the whole while it is outgrown in each piece. The wash moves the minus, it does not delete it.

\begin{table}[ht]
\centering
\begin{tabular}{@{}llll@{}}
\toprule
reading & duration & exit & overall \\
\midrule
purgative release & time to encode the turn & the lean lifts you when the wash completes & highest \\
self-chosen clinging & as long as the self clings & releasing the clinging & high \\
recurrence loop & many loops, bounded by the leak & the cycle broken by the lean or the wake & high \\
persistent attractor & long but bounded & eventual, it cannot hold against the lean & medium \\
eternal hell & forever & never & ruled out \\
\bottomrule
\end{tabular}
\caption{How long a hell can hold, ranked.}
\label{tab:hell-duration}
\end{table}

The model's clearest answer is purgative release. Hell lasts exactly as long as the soul-washing takes, and the lean lifts you out when it is done. The harder readings explain why a hell can feel endless, but none is eternal.

\begin{principle}[The gate of hell]
The gate of hell is locked from the inside, and the lean is the current that carries you across the finish line.
\end{principle}

\subsection{Heaven, the lean's home}

Heaven mirrors hell. It is many, dynamic, and conserved.

A heaven is a pleasure-region, a bright cloud of $+1$ tone. The lean pulls toward the warm tones, so a self drifts toward the bright clouds. It does not have to be turned into them the way it must be turned into the dark ones.

Heaven is not static. The model has no eternal stasis, so even a heaven is passed through, not frozen in. A self drawn into a bright cloud is soothed, but the dynamics keeps moving and the mesh keeps growing at its wake, so the self is eventually carried on. This is gentler than it sounds. The path rises, so being carried on from a heaven is usually being carried toward a better one.

The highest heaven is the deep aligned summit, the region high on the integration ladder and warm in tone, where a self is both maximally integrated and at rest in the lean. It is the nearest the model comes to the supreme paradise, union and bliss together. It sits exactly where the two verticals meet the warm tone.

And heaven is paid for. The total pleasure exactly equals the total pain, so every heaven is balanced by a hell elsewhere. There is no universal heaven. The lean can draw you into a bright cloud, but it cannot raise the books.

\subsection{The perfection spiral}

A self does not sit in one realm. It moves, a spiral through the clouds and up the ladder.

Track it by three coordinates. Its tone, pain or peace or pleasure, oscillates. Its integration, its level on the tower, oscillates. And its alignment, how well its navigation flows with the lean, is the perfection coordinate.

The motion has two parts. The first is the oscillation. Tone and integration swing. The self rises into a heaven, falls into a hell, is washed, rises again. This is the wheel, the cyclic part, the reversible dynamics turning the self through states. The second is the drift. The lean biases the dynamics, so over many oscillations the alignment trends upward. Each wash encodes a lesson and re-routes the self a little more toward the good.

The composition of an oscillation with an upward drift is a spiral. The self circles through the states, up the ladder and down, but each loop closes a little higher in alignment than the last. The wheel of experience climbs as it turns.

\begin{result}[The descents are the climb]
The descents are not regressions. They are the mechanism of the climb. A self falls into a hell, the wash encodes a lesson, and it comes out higher than it went in. The good direction runs through the dark clouds, not around them.
\end{result}

Crucially, progress is organization, not pleasure. Conservation forbids net pleasure. So a self perfecting is not getting richer in the warm tone. It is getting clearer, more unified, more aligned, wiser. \textbf{Purity is the absence of misaligned structure}, and the soul-washing is exactly its removal. The self becomes a better vessel of the conserved field, not a hoarder of its pluses.

The self approaches perfection forever and never finally sits in it. There is no eternal stasis and no complete self-model, and the world grows forever so the target itself recedes. The journey is the perfection, an inexhaustible ascent toward a horizon that forever recedes.

\subsection{The honest bottom line}

The solid core is unusually strong here. Conservation balancing to peace, the temporariness of hell, the multiplicity of clouds, the purgative lean, and the upward spiral are forced by the model, not chosen. These are the committed tier.

The plausible middle rests on persistence, the lessons must last. The speculative edge is the felt redemptive quality of the wash, riding on the experiential posit. The ruled-out region is exactly the cruel and the magical. An eternal hell. A free heaven that raises the books. Pure punishment with no transform. A self sent in from outside.

\begin{principle}[The refined myth]
The model's heaven and hell are temporary, many, conserved, purgative, and spiraling toward an endless perfection. On the geometry's own terms, the picture is both more honest and more merciful than the myth it refines.
\end{principle}
